\section{Experiment RMIG-FBOHN v1}
\label{sec:rmig_fbohn_v1}

The first benchmark used a label constructed from orbital energies and Laplacian energies. Mean results:
\begin{center}
\begin{tabular}{lrr}
\toprule
Model & Accuracy & Interpretation\\
\midrule
RAW & 0.491 & no access to orbital structure\\
BOHN & 0.803 & orbital energy is informative\\
FBOHN & 0.989 & RMIG Laplacian geometry almost completely explains the label\\
FBOHN-SRL-12 & 0.871 & PCA compresses some but not all information\\
\bottomrule
\end{tabular}
\end{center}

\subsection{Conclusion}

The first test confirmed that FBOHN is not merely an enlargement of BOHN. Adding Laplacian channels $S_O$, $L_O$ contributes new geometric information. The jump
\[
0.491\to 0.803\to 0.989
\]
is, however, partially expected because the label was constructed from FBOHN features. Therefore a methodologically harder test was performed.
